%=============================================================================
\section{Discussion}
\label{sec:discussion}
%=============================================================================

\subsection{What Spectral Metrics Tell You That Loss Cannot}

\begin{table}[t]
\centering
\small
\caption{Capabilities comparison: loss-only monitoring vs.\ spectral metrics (SR/$d$ + $\alpha$).}
\label{tab:compare}
\begin{tabular}{@{}p{4.5cm}cc@{}}
\toprule
\textbf{Capability} & \textbf{Loss} & \textbf{SR/$d$ + $\alpha$} \\
\midrule
Detect training progress & \checkmark & \checkmark \\
Detect structural learning & $\times$ & \checkmark \\
Detect structural \emph{degradation} & $\times$ & \checkmark ($\alpha$ reversal) \\
Adaptive schedule trigger & $\times$ & \checkmark ($d\alpha/dt = 0$) \\
Predict downstream perf.\ without eval & $\times$ & \checkmark ($r = -0.92$) \\
Compare across model sizes & $\times$ & \checkmark (normalized) \\
Diagnose layer-specific issues & $\times$ & \checkmark ($\alpha_{\mathrm{attn}}$ vs.\ $\alpha_{\mathrm{mlp}}$) \\
Determine training sufficiency & $\times$ & \checkmark ($\alpha < 3$) \\
\bottomrule
\end{tabular}
\end{table}

The fundamental asymmetry is: \emph{loss can decrease while structure degrades}. This occurs whenever the model reduces prediction error through memorization rather than compression. Only spectral metrics distinguish these cases.

\subsection{Connection to Over-Training and Chinchilla}

The Chinchilla scaling law~\citep{hoffmann2022} minimizes loss for a given compute budget, yielding $D/N \approx 20$ tokens/parameter. Our Structural Chinchilla law (Eq.~\ref{eq:structural_chinchilla}) reveals that this produces $\alpha \approx 5.8$---models that are \emph{compute-optimal but structurally immature}.

The industry has empirically moved toward massive over-training: Llama-3 uses $D/N \approx 15{,}000$ (750$\times$ Chinchilla), and these models exhibit superior generalization, fine-tuning robustness, and emergent capabilities~\citep{llama3}. Our framework explains why: only at $D/N > 300$ does the spectral structure enter the heavy-tail regime ($\alpha < 4$), which corresponds to robust, transferable feature representations.

This resolves the apparent paradox of ``compute-inefficient'' training being better: \textbf{Chinchilla minimizes loss per FLOP, but structural optimality requires a fundamentally different budget.}

\subsection{Practical Monitoring Protocol}

Based on our findings, we recommend a three-signal protocol for training teams:

\begin{enumerate}[leftmargin=*,itemsep=2pt]
    \item \textbf{$\alpha$ (primary, every 5K steps)}: Track global $\alpha$ and layer-type decomposition.
    \begin{itemize}[leftmargin=*,itemsep=0pt]
        \item $d\alpha/dt < 0$: structure forming (healthy)
        \item $d\alpha/dt \approx 0$: structure saturated (consider decay)
        \item $d\alpha/dt > 0$ for 3+ measurements: \textbf{ALERT}---structural degradation
        \item $\alpha_{\mathrm{mlp}} \gg \alpha_{\mathrm{attn}}$: MLP is the capacity bottleneck
    \end{itemize}

    \item \textbf{SR/$d$ (secondary, every 10K steps)}: Monitor compression progress.
    \begin{itemize}[leftmargin=*,itemsep=0pt]
        \item SR/$d$ still decreasing: useful compression ongoing
        \item SR/$d$ plateaued at ${\sim}0.05$: maximum compression achieved
        \item SR/$d$ < 0.03: potential over-compression (reduce weight decay)
    \end{itemize}

    \item \textbf{Performance prediction (on demand)}:
    Use Eq.~\ref{eq:regression} to estimate downstream performance from current SR/$d$ and $N$, without running evaluation.
\end{enumerate}

\subsection{Limitations}

\paragraph{Scale ceiling.} Our measurements extend to 7B parameters. Whether SR/$d \approx 0.055$ holds at 70B+ remains unverified, though the asymptotic model (SR/$d = 0.04 + 0.61/\sqrt{d}$) predicts convergence to 0.04 for very large $d$.

\paragraph{Architecture scope.} We validate across GPT-NeoX, LLaMA, and OLMo2---all dense decoder-only transformers. MoE and encoder-decoder architectures are untested.

\paragraph{Proxy training.} The $\alpha$-guided experiment (\S\ref{sec:schedule}) uses random tokens as a proxy. A definitive validation requires real data with downstream evaluation, which we plan as future work.

\paragraph{Correlation vs.\ causation.} SR/$d$ correlates with but does not provably cause good downstream performance. The causal mechanism---whether compression is a \emph{driver} or \emph{side effect} of learning---deserves further theoretical analysis.

\paragraph{Power-law fitting.} The $\alpha$ exponent depends on the fitting procedure (choice of $x_{\min}$, tail range). We use the MLE + KS approach from~\citet{martin_mahoney_jmlr2021}, but report that simpler log-log regression produces qualitatively similar trends with slightly different absolute values.
